\section{Conventions}\label{sec:conventions}

All varieties are over $\CC$ and all cohomology is singular cohomology with
the indicated coefficients. For a smooth projective variety $X$ we write
$H^{p,q}(X)$ for the Hodge components of $H^{p+q}(X,\CC)$, and
\[
  \Hdg^{p}(X)=H^{2p}(X,\QQ)\cap H^{p,p}(X)
\]
for the space of rational Hodge classes of type $(p,p)$. A class is
\emph{algebraic} if it lies in the image of the cycle class map
$\cl^{p}_{X}\colon\CH^{p}(X)_{\QQ}\to\Hdg^{p}(X)$; the subspace of algebraic
classes is a subring of $H^{*}(X,\QQ)$, since the cycle class map is a ring
homomorphism on the level of intersection products.%
\footnote{The intersection product on $\CH^{*}(X)_{\QQ}$ exists because $X$ is
smooth, and the compatibility of the cycle class map with it is
\cite[Chapter~19]{Ful98}. The integral form of the question, with $\ZZ$ in
place of $\QQ$, is a different problem and is not considered anywhere in this
paper.}

An abelian variety $A$ of dimension $g$ has $H^{1}(A,\QQ)$ of dimension $2g$
and a canonical isomorphism
\begin{equation}\label{eq:wedge}
  H^{k}(A,\QQ)\;\cong\;\textstyle\bigwedge^{k}H^{1}(A,\QQ)
\end{equation}
of Hodge structures, with the product on the left corresponding to the wedge
product on the right. We use \eqref{eq:wedge} without further comment; it is
the single structural fact that makes every computation below finite.%
\footnote{Topologically $A=V/\Lambda$ is a real torus of dimension $2g$, so
$H^{*}(A,\ZZ)=\bigwedge^{*}\Hom(\Lambda,\ZZ)$ by the K\"unneth formula for
$(S^{1})^{2g}$, with the cup product corresponding to the wedge product. The
Hodge decomposition of $H^{k}$ is the $k$-th exterior power of
$H^{1}=H^{1,0}\oplus H^{0,1}$, because the translation invariant forms
$dz_{I}\wedge d\bbar z_{J}$ are harmonic for the flat metric and represent
every class.}
The dual abelian variety is $\Av=\Pic^{0}(A)$, and
$H^{1}(\Av,\QQ)\cong H^{1}(A,\QQ)^{\vee}$.

A polarisation of $A$ is a class $\theta\in\NS(A)\otimes\QQ$ that is ample;
equivalently an alternating form $E$ on $H_{1}(A,\QQ)$ satisfying the Riemann
relations. The polarisation is principal when $E$ is unimodular on
$H_{1}(A,\ZZ)$, and then $\theta^{g}/g!$ is the class of a point. We write
$\theta'$ for the polarisation on $\Av$ induced by $\theta$ through the
isogeny $\phi_{\theta}\colon A\to\Av$; when $\theta$ is principal,
$\phi_{\theta}$ is an isomorphism and $\phi_{\theta}^{*}\theta'=\theta$.

For an endomorphism $f$ of $A$ we write $f^{*}$ for the induced map on
cohomology. When $K$ is a field acting on $A$ through
$\iota\colon K\hookrightarrow\End^{0}(A)=\End(A)\otimes\QQ$, the element
$\iota(\tau)$ for $\tau\in K$ is an element of $\End^{0}(A)$; it acts on
$H^{*}(A,\QQ)$, and we allow ourselves to speak of the action of $\tau$ when
$\iota$ is understood. For $\tau$ in the order $\iota^{-1}(\End(A))$ the
element $\iota(\tau)$ is an honest isogeny, and the general case follows by
$\QQ$-linearity; no statement below depends on the distinction.

The Chern character of a bounded complex of coherent sheaves is
$\ch(\cE^{\bullet})=\sum_{i}(-1)^{i}\ch(\cE^{i})$ for any bounded locally free
resolution, and this is well defined on $D^{b}$; on a smooth projective
variety every coherent sheaf admits such a resolution. We write $\ch_{k}$ for
the component in $H^{2k}$.%
\footnote{On an abelian variety the tangent bundle is trivial and so is the
Todd class. Grothendieck--Riemann--Roch therefore reads
$\ch(f_{!}x)=f_{*}\ch(x)$ for a morphism $f$ between abelian varieties, and
$\ch\colon K_{0}(A)_{\QQ}\to\CH^{*}(A)_{\QQ}$ is an isomorphism of rings
\cite[Chapter~15]{Ful98}. Both facts are used below without further mention;
for the inclusion of a smooth subvariety, whose normal bundle need not be
trivial, the Todd class of the normal bundle does enter, and it is written out
where it matters.}

For the Fourier--Mukai transform we follow \cite{Muk81}. The Poincar\'e bundle
$\cP$ on $A\times\Av$ is normalised so that $\cP|_{A\times\{0\}}$ and
$\cP|_{\{0\}\times\Av}$ are trivial. With $p\colon A\times\Av\to A$ and
$q\colon A\times\Av\to\Av$ the projections, the transform is
\[
  \Phi_{\cP}\colon D^{b}(\Av)\to D^{b}(A),
  \qquad
  \Phi_{\cP}(\cG)=Rp_{*}\bigl(Lq^{*}\cG\otimes^{L}\cP\bigr).
\]
On cohomology we write the induced map with the same symbol,
\[
  \Phi_{\cP}\colon H^{*}(\Av,\QQ)\longrightarrow H^{*}(A,\QQ),
  \qquad
  \Phi_{\cP}(x)=p_{*}\bigl(q^{*}x\cdot\ch(\cP)\bigr),
\]
and the Grothendieck--Riemann--Roch theorem gives
$\ch(\Phi_{\cP}(\cG))=\Phi_{\cP}(\ch(\cG))$ because the Todd class of an
abelian variety is trivial; see \Cref{prop:grr}.

We write $\Nm=\Nm_{K/\QQ}$ for the norm of an imaginary quadratic field and
$\bbar{\tau}$ for the conjugate of $\tau\in K$. Characters of $\Kx$ are written
multiplicatively; $\Nm^{n}$ denotes $\tau\mapsto\Nm(\tau)^{n}$.

Finally, a word on what we mean by a construction being \emph{of ambient type}.
We mean precisely the hypothesis of \Cref{thm:divisor}: a morphism
$i\colon\Av\to\PP^{N-1}$ defined by a linear system, an object $\cF$ of
$D^{b}(\PP^{N-1})$, and the class $\ch_{k}(\Phi_{\cP}(Li^{*}\cF))$. Nothing is
assumed about how $\cF$ is chosen, and in particular $\cF$ may be the
structure sheaf of a secant variety, of a join, of a determinantal locus, or
of any other subscheme, and it may be twisted by any line bundle from the
ambient space.
